\usepackage[T1]{fontenc}
\usepackage[utf8]{inputenc}
\usepackage{lmodern}
\usepackage{amsmath,amssymb,amsthm,mathtools}
\usepackage[letterpaper,margin=1.08in]{geometry}
\usepackage{microtype}
\usepackage{enumitem}
\usepackage{booktabs,array,tabularx}
\usepackage{xcolor}
\usepackage{hyperref}
\usepackage{bookmark}
\input{glyphtounicode}
\pdfglyphtounicode{mu}{03BC}
\pdfglyphtounicode{Delta}{0394}
\pdfglyphtounicode{backslash}{2216}
\pdfglyphtounicode{mapsto}{007C}
\pdfglyphtounicode{parenleftbig}{0028}
\pdfglyphtounicode{parenrightbig}{0029}
\pdfglyphtounicode{braceleftbig}{007B}
\pdfglyphtounicode{bracerightbig}{007D}
\pdfglyphtounicode{parenleftbigg}{0028}
\pdfglyphtounicode{parenrightbigg}{0029}
\pdfglyphtounicode{braceleftbigg}{007B}
\pdfglyphtounicode{bracerightbigg}{007D}
\pdfglyphtounicode{parenleftBigg}{0028}
\pdfglyphtounicode{parenrightBigg}{0029}
\pdfglyphtounicode{braceleftBigg}{007B}
\pdfglyphtounicode{bracerightBigg}{007D}
\pdfglyphtounicode{vextendsingle}{23D0}
\pdfglyphtounicode{vextenddouble}{2016}
\pdfglyphtounicode{summationtext}{2211}
\pdfglyphtounicode{summationdisplay}{2211}
\pdfglyphtounicode{logicalordisplay}{22C1}
\pdfglyphtounicode{bracketleftBig}{005B}
\pdfglyphtounicode{bracketrightBig}{005D}
\pdfglyphtounicode{bracketleftbig}{005B}
\pdfglyphtounicode{bracketrightbig}{005D}
\pdfglyphtounicode{bracketleftBigg}{005B}
\pdfglyphtounicode{bracketrightBigg}{005D}
\pdfglyphtounicode{integraltext}{222B}
\pdfglyphtounicode{integraldisplay}{222B}
\pdfglyphtounicode{radicalbig}{221A}
\pdfglyphtounicode{radicalBig}{221A}
\pdfglyphtounicode{bracelefttp}{23A7}
\pdfglyphtounicode{braceleftmid}{23A8}
\pdfglyphtounicode{braceleftbt}{23A9}
\pdfglyphtounicode{braceex}{23AA}
\pdfgentounicode=1
\hypersetup{
  colorlinks=true,
  linkcolor=blue!45!black,
  citecolor=blue!45!black,
  urlcolor=blue!45!black,
  pdftitle={Sharp nodal length on smooth surfaces},
  pdfauthor={OpenAI},
  pdfsubject={Laplace eigenfunctions, subharmonic profiles, and nodal length},
  pdfcreator={pdfLaTeX},
  bookmarksnumbered=true,
  pdfdisplaydoctitle=true
}
\newcommand{\R}{\mathbb R}
\newcommand{\N}{\mathbb N}
\newcommand{\D}{\mathbb D}
\newcommand{\E}{\mathbb E}
\newcommand{\eps}{\varepsilon}
\DeclareMathOperator{\av}{avg}
\DeclareMathOperator{\supp}{supp}
\DeclareMathOperator{\dist}{dist}
\DeclarePairedDelimiter{\norm}{\lVert}{\rVert}
\DeclarePairedDelimiter{\abs}{\lvert}{\rvert}
\newtheorem{theorem}{Theorem}[section]
\newtheorem{proposition}[theorem]{Proposition}
\newtheorem{lemma}[theorem]{Lemma}
\newtheorem{corollary}[theorem]{Corollary}
\newtheorem{claim}[theorem]{Claim}
\theoremstyle{definition}
\newtheorem{definition}[theorem]{Definition}
\theoremstyle{remark}
\newtheorem{remark}[theorem]{Remark}
\numberwithin{equation}{section}
\setlist[enumerate]{itemsep=3pt,topsep=6pt}
\setlist[itemize]{itemsep=3pt,topsep=6pt}
\setlength{\emergencystretch}{1.5em}
\allowdisplaybreaks[1]
\widowpenalty=8000
\clubpenalty=8000
\makeatletter
\renewcommand{\@maketitle}{%
  \newpage
  \null
  \vskip 2em
  \begin{center}%
    {\LARGE \@title\par}%
    \vskip 0.75em
    {\large \@author\par}%
    \vskip 0.35em
    {\normalsize \@date\par}%
  \end{center}%
  \vskip 1em
}
\makeatother
